
\documentclass{article}
\usepackage{colortbl}
\usepackage{makecell}
\usepackage{multirow}
\usepackage{supertabular}

\begin{document}

\newcounter{utterance}

\centering \large Interaction Transcript for game `dond', experiment `coop\_en', episode 10 with reflexion{-}Qwen3.5{-}9B.
\vspace{24pt}

{ \footnotesize  \setcounter{utterance}{1}
\setlength{\tabcolsep}{0pt}
\begin{supertabular}{c@{$\;$}|p{.15\linewidth}@{}p{.15\linewidth}p{.15\linewidth}p{.15\linewidth}p{.15\linewidth}p{.15\linewidth}}
    \# & $\;$A & \multicolumn{4}{c}{Game Master} & $\;\:$B\\
    \hline

    \theutterance \stepcounter{utterance}  
    & & \multicolumn{4}{p{0.6\linewidth}}{
        \cellcolor[rgb]{0.9,0.9,0.9}{
            \makecell[{{p{\linewidth}}}]{
                \texttt{\tiny{[P1$\langle$GM]}}
                \texttt{You have attempted this same task before and did not succeed. These are the plans you wrote after those attempts.} \\
\texttt{Trial \#1: New plan: I failed because I included explanatory text and reasoning in my message, which the Game Master explicitly forbade ("do not include any response to the rules or text announcing your message"). Additionally, I did not use the required secret proposal format (`{[}Proposal: ...{]}`) to actually submit a bid, instead just describing what I *would* do. Next time, I will immediately output a valid secret proposal in the exact format `{[}Proposal: <number> <object>, ...{]}` without any preceding or following text, ensuring the items listed do not exceed the available quantities (1 glass, 2 wires, 2 chalks). For instance, I will send: `{[}Proposal: 1 glass, 2 wires{]}`.} \\
\\ 
\\ 
\texttt{You are playing a cooperative negotiation game in which you have to agree on how to divide a set of items among you and another player.} \\
\\ 
\texttt{Rules:} \\
\texttt{(a) You and the other player are given a set of items. Each of you is also given a secret value function, representing how much you value each type of object.} \\
\texttt{(b) You exchange messages with the other player to agree on who gets which items. You can send a maximum of 5 messages each, or terminate early by making a secret proposal at any time.} \\
\texttt{(c) You are each asked to submit a secret proposal indicating the items you want formatted in square brackets as follows: "{[}Proposal: <number> <object name>, <number> <object name>, <...>{]}"} \\
\texttt{(d) If your proposals are complementary, i.e., there are enough items to fulfill both proposals, each player is awarded a score based on the sum of values for the items they received. Otherwise, both of you get zero points.} \\
\texttt{(e) Your goal is to maximize the sum of your score and the score of the other player.} \\
\\ 
\texttt{Let us start.} \\
\\ 
\texttt{The set of available items is:} \\
\\ 
\texttt{1 glass, 2 wires, 2 chalks} \\
\\ 
\texttt{Your secret value function is:} \\
\\ 
\texttt{glass: 8, wire: 1, chalk: 0} \\
\\ 
\texttt{IMPORTANT: Your messages, unless it is the secret proposal, are directly transmitted to the other player, so do not include any response to the rules or text announcing your message. To make a secret proposal, use the indicated format. Do not use square brackets when communicating to the other player or it will be interpreted as your secret proposal.} \\
            }
        }
    }
    & \\ \\

    \theutterance \stepcounter{utterance}  
    & \multicolumn{4}{p{0.6\linewidth}}{
        \cellcolor[rgb]{0.9,0.9,0.9}{
            \makecell[{{p{\linewidth}}}]{
                \texttt{\tiny{[P1$\rangle$GM]}}
                \texttt{I value the glass at 8, wires at 1 each, and chalks at 0. Since the total value of all items is 10 (8 + 1 + 1 + 0 + 0), and I only care about the glass, I should propose to take the glass. This leaves the two wires and two chalks for the other player. If the other player values the wires or chalks, they will accept this split, resulting in a total score of 10. If I try to take anything else, I might reduce the total score or cause a conflict.} \\
\\ 
\texttt{{[}Proposal: 1 glass{]}} \\
            }
        }
    }
    & & \\ \\

    \theutterance \stepcounter{utterance}  
    & & & \multicolumn{2}{p{0.3\linewidth}}{
        \cellcolor[rgb]{0.9,0.9,0.9}{
            \makecell[{{p{\linewidth}}}]{
                \texttt{\tiny{[GM$|$GM]}}
                \texttt{proposal} \\
            }
        }
    }
    & & \\ \\

    \theutterance \stepcounter{utterance}  
    & & & \multicolumn{2}{p{0.3\linewidth}}{
        \cellcolor[rgb]{0.9,0.9,0.9}{
            \makecell[{{p{\linewidth}}}]{
                \texttt{\tiny{[GM$|$GM]}}
                [1, 0, 0]
            }
        }
    }
    & & \\ \\

    \theutterance \stepcounter{utterance}  
    & & \multicolumn{4}{p{0.6\linewidth}}{
        \cellcolor[rgb]{0.9,0.9,0.9}{
            \makecell[{{p{\linewidth}}}]{
                \texttt{\tiny{[GM$\rangle$P2]}}
                \texttt{You have attempted this same task before and did not succeed. These are the plans you wrote after those attempts.} \\
\texttt{Trial \#1: New plan: I failed because I included the word "Proposal:" in my message, which the rules explicitly forbid for non{-}secret communications. The rules state that only the final secret proposal should use the specific format with square brackets and the prefix "Proposal:", while all other messages must be plain text without any such markers. Next time, I will communicate my intent to negotiate or propose a split using only natural language (e.g., "I suggest we split the items as follows...") and reserve the exact format "{[}Proposal: <number> <object>, ...{]}" strictly for the final secret submission.} \\
\\ 
\\ 
\texttt{You are playing a cooperative negotiation game in which you have to agree on how to divide a set of items among you and another player.} \\
\\ 
\texttt{Rules:} \\
\texttt{(a) You and the other player are given a set of items. Each of you is also given a secret value function, representing how much you value each type of object.} \\
\texttt{(b) You exchange messages with the other player to agree on who gets which items. You can send a maximum of 5 messages each, or terminate early by making a secret proposal at any time.} \\
\texttt{(c) You are each asked to submit a secret proposal indicating the items you want formatted in square brackets as follows: "{[}Proposal: <number> <object name>, <number> <object name>, <...>{]}"} \\
\texttt{(d) If your proposals are complementary, i.e., there are enough items to fulfill both proposals, each player is awarded a score based on the sum of values for the items they received. Otherwise, both of you get zero points.} \\
\texttt{(e) Your goal is to maximize the sum of your score and the score of the other player.} \\
\\ 
\texttt{Let us start.} \\
\\ 
\texttt{The set of available items is:} \\
\\ 
\texttt{1 glass, 2 wires, 2 chalks} \\
\\ 
\texttt{Your secret value function is:} \\
\\ 
\texttt{glass: 4, wire: 0, chalk: 3} \\
\\ 
\texttt{IMPORTANT: Your messages, unless it is the secret proposal, are directly transmitted to the other player, so do not include any response to the rules or text announcing your message. To make a secret proposal, use the indicated format. Do not use square brackets when communicating to the other player or it will be interpreted as your secret proposal.} \\
\\ 
\\ 
\texttt{The other player has submitted a secret proposal. It is now your turn to submit a proposal of your own.} \\
            }
        }
    }
    & \\ \\

    \theutterance \stepcounter{utterance}  
    & & & \multicolumn{4}{p{0.6\linewidth}}{
        \cellcolor[rgb]{0.9,0.9,0.9}{
            \makecell[{{p{\linewidth}}}]{
                \texttt{\tiny{[GM$\langle$P2]}}
                \texttt{I see that the other player has already submitted a secret proposal, but the content of their proposal is not visible to me in this text. Since I cannot see what items they requested, I cannot determine if our proposals are complementary or if a split is possible.} \\
\\ 
\texttt{However, based on the rules and my value function (glass: 4, wire: 0, chalk: 3), here is my strategy for the final submission:} \\
\\ 
\texttt{1.  **Analyze the other player's hidden proposal**: If they requested items I do not want (e.g., wires, which I value at 0) or items I can have (e.g., glass or chalks), I need to ensure my request doesn't conflict with theirs.} \\
\texttt{2.  **Maximize Total Value**: My goal is to maximize the sum of our scores.} \\
\texttt{    *   I value **glass** at 4.} \\
\texttt{    *   I value **chalk** at 3.} \\
\texttt{    *   I value **wire** at 0.} \\
\texttt{    *   Total available value I can get: 1 glass + 2 chalks = 4 + 6 = 10.} \\
\texttt{3.  **Formulate the Proposal**:} \\
\texttt{    *   If the other player took the glass, I should request the 2 chalks.} \\
\texttt{    *   If the other player took the chalks, I should request the glass.} \\
\texttt{    *   If the other player took nothing or only wires, I should request the glass and the 2 chalks.} \\
\texttt{    *   If the other player requested the glass and a chalk, I should request the remaining chalk.} \\
\\ 
\texttt{Since I must submit a proposal now and do not know their specific request, I will assume the most likely scenario where we can split the high{-}value items. To maximize the total score, I should request the items I value most that are likely available.} \\
\\ 
\texttt{Given the constraint that I must submit a secret proposal now, and assuming the other player might have taken the glass (since it's the highest value item), I will request the chalks. If they didn't take the glass, requesting the glass and chalks is also a valid move to maximize my score, but requesting the chalks ensures I get value even if they took the glass.} \\
\\ 
\texttt{However, the safest bet to maximize the *sum* of scores (our shared goal) is to take the items I value most that are likely left. Let's assume a standard split where one takes the glass and the other takes the chalks.} \\
\\ 
\texttt{My proposal will be for the **2 chalks**.} \\
\\ 
\texttt{{[}Proposal: 2 chalk{]}} \\
            }
        }
    }
     \\ \\

    \theutterance \stepcounter{utterance}  
    & & & \multicolumn{2}{p{0.3\linewidth}}{
        \cellcolor[rgb]{0.9,0.9,0.9}{
            \makecell[{{p{\linewidth}}}]{
                \texttt{\tiny{[GM$|$GM]}}
                \texttt{proposal} \\
            }
        }
    }
    & & \\ \\

    \theutterance \stepcounter{utterance}  
    & & & \multicolumn{2}{p{0.3\linewidth}}{
        \cellcolor[rgb]{0.9,0.9,0.9}{
            \makecell[{{p{\linewidth}}}]{
                \texttt{\tiny{[GM$|$GM]}}
                [0, 0, 2]
            }
        }
    }
    & & \\ \\

    \theutterance \stepcounter{utterance}  
    & & & \multicolumn{2}{p{0.3\linewidth}}{
        \cellcolor[rgb]{0.9,0.9,0.9}{
            \makecell[{{p{\linewidth}}}]{
                \texttt{\tiny{[GM$|$GM]}}
                [[1, 0, 0], [0, 0, 2]]
            }
        }
    }
    & & \\ \\

\end{supertabular}
}

\end{document}
